\documentclass{abntex2}
\usepackage{fontspec}
\setmainfont{Times New Roman}[
  BoldFont      = Times New Roman Bold,
  ItalicFont    = Times New Roman Italic,
  BoldItalicFont= Times New Roman Bold Italic
]
\usepackage{microtype}
\usepackage[a4paper,top=3cm,bottom=2cm,left=3cm,right=2cm]{geometry}

% Fontes (usa variável do Pandoc ou fallback)
\usepackage{fontspec}
$if(mainfont)$
  \setmainfont{$mainfont$}
$else$
  \setmainfont{Times New Roman}
$endif$

\usepackage{microtype}
\usepackage{polyglossia}
\setdefaultlanguage{brazil}

% Metadados (Pandoc)
\titulo{$title$}
\autor{$author$}
\orientador{$advisor$}
\instituicao{$institution$}
\local{$place$}
\data{$year$}

% Instituição padrão se vazia
\makeatletter
\@ifundefined{ABNT@instituicao}{\instituicao{Universidade de São Paulo}}{}
\makeatother

\begin{document}
\imprimircapa
\imprimirfolhaderosto*

% Sumário / listas
\pdfbookmark[0]{Sumário}{sumario}
\tableofcontents*
$if(lof)$\listoffigures*$endif$
$if(lot)$\listoftables*$endif$

% Corpo
$body$

% Referências
$if(bibliography)$
\bibliographystyle{abntex2-alf}
\bibliography{$for(bibliography)$$bibliography$$sep$,$endfor$}
$endif$

\end{document}
% ==== CAPA PADRONIZADA USP/ABNT COM ORIENTADOR ====
\makeatletter
\renewcommand{\imprimircapa}{%
  \begin{capa}
    \centering
    {\fontsize{14}{16}\selectfont\textbf{\ABNTinstituicao}\par}
    \vspace*{2.5cm}
    {\fontsize{12}{14}\selectfont\ABNTautor\par}
    \vfill
    {\fontsize{18}{20}\selectfont\textbf{\ABNTtitulo}\par}
    \vfill
    {\fontsize{12}{14}\selectfont Orientador: \@orientador\par}
    \vfill
    {\fontsize{12}{14}\selectfont\ABNTlocal\par}
    {\fontsize{12}{14}\selectfont\ABNTdata\par}
  \end{capa}%
}
\makeatother
% ==== AJUSTES DO SUMÁRIO (ABNT 12 pt) ====
\makeatletter
% fonte 12 pt para níveis do sumário
\renewcommand{\cftsecfont}{\normalsize}
\renewcommand{\cftsecpagefont}{\normalsize}
\renewcommand{\cftsubsecfont}{\normalsize}
\renewcommand{\cftsubsecpagefont}{\normalsize}
% pontilhado mais “cheio”
\renewcommand{\cftdotsep}{1}
% reduzir espaçamento vertical entre entradas
\setlength{\cftbeforesecskip}{0pt}
\setlength{\cftbeforesubsecskip}{0pt}
\makeatother
% --- Ajustes finos do Sumário (abnTeX2) ---
\makeatletter
\renewcommand{\cftchapterfont}{\normalsize}
\renewcommand{\cftsectionfont}{\normalsize}
\renewcommand{\cftsubsectionfont}{\normalsize}
\renewcommand{\cftdotsep}{1}
\renewcommand{\cftchapindent}{0pt}
\renewcommand{\cftsecindent}{0pt}
\renewcommand{\cftsubsecindent}{0pt}
\setlength{\cftbeforechapskip}{6pt}
\setlength{\cftbeforesecskip}{2pt}
\setlength{\cftbeforesubsecskip}{1pt}
\makeatother
